\makeabstract
{
AI Going Against the Grain: Testing a Neural Network’s Ability to Identify Microtextures in Quartz Sand 
}
{
Drew Weisgerber (1), David Jones (1), Hanna Brooks (1)
}
{
\textsuperscript{1}Department of Geology, Amherst College, Amherst, MA, USA
}
{
As grains of sand are transported across Earth’s surface, they acquire distinctive surface textures characteristic of their environment. Rivers, glaciers, waves, and wind all produce different textures. The distribution of microtextures on individual grains can be used to reconstruct the sedimentary history of a sediment. SAND AI is a Deep Neural Network (developed by Hasson et al. 2024) intended to address the limitations of microtextural analysis by providing an objective and qualitative method of determining depositional environments. This study investigated the extent to which successive sedimentary cycles can be derived from the model’s output, using samples with known sedimentary histories.
}
{
The model was developed and trained (by Hasson et al. 2024) to distinguish between glacial, eolian, fluvial, and beach transport, using sand grains from known environments. It then underwent external validation on sediment from modern sites unknown to it. We used SAND AI to analyze sediments from environments it was trained to recognize, in addition to those it was not. Each sample was cleaned in hydrogen peroxide and boiling HCl in order to remove surface coatings. Grains were imaged both before and after cleaning in order to gauge the effectiveness of the cleaning procedure. Quartz grains were randomly selected
under a microscope and imaged on a Scanning Electron Microscope (SEM). A spectrum was taken on a clean face of every grain analyzed to verify that it was quartz. Images were then processed by the SAND AI model. 
}
{
The model proved better at accurately identifying glacial and eolian environments than beach and fluvial ones, and we found that glacial microtextures linger on grains that are subsequently transported in other environments. Samples for which fewer than 5 grains were identified at >75\% confidence are considered to have not been successfully identified by the model. The one lithified sample tested was incorrectly identified, which may be due to diagenetic alteration of microtextures and cement growth. 
}
{
Though we conclude that SAND AI can determine a sediment’s past environments, with the current model output we cannot confidently form a chronological sedimentary history without further observation, because it does not recognize cross-cutting relationships between microtextures. In addition, more work is needed to determine if the model can successfully be applied to uncover the transport history of lithified sedimentary rocks, or if diagenetic alteration renders the model inaccurate.
}

\newpage
